\practicechapter{第 11 章实践：整数指令、位运算和地址计算}

本章对应第一册第 11 章“整数指令、位运算和地址计算”。实践目标有两条：一是读懂反汇编里的整数加法、位移、按位与或和地址计算；二是回到解析器源码，理解小端整数读取本身就是一组位运算。

\practicesection{一、对应原书问题：整数不是只有加减乘除}

ELF 解析器中到处都是整数：offset、size、address、flags、type、binding、section index。它们有的表示数量，有的表示地址，有的是位标志，有的是枚举值。若把所有整数都当普通十进制数看，就看不出格式结构。

源码中的 \code{read\_u32\_le} 展示了最基础的位运算：取一个字节，转换成更宽整数，再左移 \code{i * 8} 位，最后按位或到结果里。符号解析中的 \code{info \& 0x0F} 和 \code{info >> 4} 则展示了同一个字节里拆出 type 和 binding 的方法。

\practicesection{二、从特例推到规律：一个字节可以装两个字段}

特例：ELF symbol 的 \code{st\_info} 字节高 4 位表示 binding，低 4 位表示 type。若 \code{info} 是 \code{0x12}，低 4 位是 \code{2}，高 4 位是 \code{1}。这不是两个字节，也不是字符串，而是一个字节里的两个小字段。

由此推出规律：二进制格式常把多个小字段压在同一个整数里。读取只是第一步，拆字段还需要 mask 和 shift。mask 用来保留关心的位，shift 用来把高位字段移动到从 0 开始的数值位置。理解这一点后，section flags、权限位、CPU flags 都会更容易读。

\practicesection{三、实践任务：解释三段整数代码}

阅读 \filepath{src/executable\_evidence.cpp} 中三段代码：

\begin{enumerate}
  \item \code{read\_u64\_le} 如何把 8 个字节合成一个 \code{uint64\_t}；
  \item \code{has\_range} 如何避免越界和整数溢出；
  \item \code{summarize\_symbols} 如何从 \code{info} 拆出 \code{type} 和 \code{binding}。
\end{enumerate}

再写一个小表，用具体数字解释 \code{info} 拆分：

\begin{lstlisting}[language=C++]
std::uint8_t info = 0x12;
std::uint8_t type = info & 0x0F;      // 2
std::uint8_t binding = info >> 4U;    // 1
\end{lstlisting}

最后在反汇编中找一条地址计算或整数操作指令，例如 \code{add}、\code{lea}、\code{and}、\code{shl}、\code{shr}。说明它的输入、输出和是否访问内存。

\practicesection{四、验收：数字要带语义}

本章交付物是 \filepath{reports/ch11-integer-instructions.md}。通过标准如下：

\begin{enumerate}
  \item 能解释 \code{read\_u64\_le} 的每次移位为什么是 8 的倍数。
  \item 能解释 \code{info \& 0x0F} 和 \code{info >> 4U} 分别保留哪部分。
  \item 能把至少一个 section flags 或 symbol type 数值翻译成人能读的含义。
  \item 能在反汇编中标注一条整数或地址计算指令。
  \item 不把所有整数都叫地址；报告要区分枚举、大小、偏移、虚拟地址和位标志。
\end{enumerate}

本章的核心训练是“数字带上下文”。没有上下文，二进制字段只是一串数；有了格式和位运算，数字才变成证据。
